\documentclass[12pt]{article}
\usepackage[margin=0.75in]{geometry}
\usepackage{fontspec}
\setmainfont{Latin Modern Roman}
\usepackage{microtype}
\usepackage{multicol}
\usepackage{fancyhdr}
\usepackage{titlesec}
\usepackage{enumitem}
\usepackage{xcolor}
\usepackage[hidelinks]{hyperref}
\setlength{\columnsep}{0.28in}
\setlength{\parindent}{1.1em}
\setlength{\parskip}{0.32em}
\setlength{\headheight}{15pt}
\titleformat{\section}{\large\bfseries\scshape}{}{0pt}{}
\titleformat{\subsection}{\normalsize\bfseries}{}{0pt}{}
\pagestyle{fancy}
\fancyhf{}
\fancyfoot[C]{\thepage}
\renewcommand{\headrulewidth}{0.4pt}
\newcommand{\focusbox}[1]{\par\vspace{0.4em}\noindent\begingroup\setlength{\fboxsep}{6pt}\colorbox{gray!12}{\begin{minipage}{\dimexpr\linewidth-2\fboxsep\relax}#1\end{minipage}}\endgroup\par\vspace{0.5em}}

\fancyhead[L]{Whately: Portfolio Revision}
\fancyhead[R]{Student Reading}
\begin{document}
\begin{center}
{\LARGE\bfseries Revising the Record of Thought}\par\vspace{0.25em}
{\large\scshape Week 18: Logic Portfolio Workshop}\par\vspace{0.5em}
{Returning to earlier work with Richard Whately's full toolkit}\par\vspace{0.6em}
\rule{0.82\textwidth}{0.4pt}
\end{center}

\section*{Central Question}
Which habits of logic have become visible in my own work?

\section*{Vocabulary Preview}
\begin{multicols}{2}
\begin{itemize}[leftmargin=*, itemsep=0.18em]
\item \textbf{Artifact}: a piece of course work that records reasoning at a particular moment.
\item \textbf{Criterion}: a standard used to judge the quality of reasoning.
\item \textbf{Diagnosis}: a precise account of where reasoning succeeds or fails.
\item \textbf{Substantive revision}: a change to thought, not merely to spelling or appearance.
\item \textbf{Reflective caption}: a short explanation of what changed, why, and what the change shows.
\item \textbf{Before-and-after evidence}: exact features of the original and revised work that make growth visible.
\item \textbf{Transfer}: using a logical habit in a new setting.
\item \textbf{Portfolio claim}: a claim about growth supported by selected artifacts.
\item \textbf{Revision trail}: original, diagnosis, revision, and caption kept together.
\item \textbf{Intellectual habit}: a repeated practice of careful judgment.
\end{itemize}
\end{multicols}

\section*{Introduction}
A portfolio is not a stack of completed pages. It is an argument about development. The evidence is your own work: an early argument skeleton, a term investigation, a syllogism map, a fallacy autopsy, an evidence test, a dispute map, a rhetoric analysis, or a Shakespeare application. Week 18 asks you to make the development visible by revising selected artifacts and explaining the revision honestly.

Revision is not the art of making an old answer look as if it had always been correct. A strong portfolio preserves the original, names a real weakness, applies a relevant Whately tool, and shows what the revised reasoning can now do.

\focusbox{\textbf{Source note:} This workshop applies the course's accumulated Whately habits to students' own artifacts. It introduces no new formal logic. The revision protocol and sample situations are original course scaffolds.}

\begin{multicols}{2}
\section*{Reading}

\subsection*{1. Choose Evidence, Not Decoration}
Do not choose an artifact only because it earned the highest score or has the neatest handwriting. Choose work that can carry a claim about growth. A useful artifact has enough visible reasoning to compare: a conclusion and assumption, an unstable term, a proposition translation, a proof map, an evidence judgment, or a literary claim.

A strong portfolio usually includes contrast. One artifact may show an early weakness; another may show later control. The goal is not perfection. The goal is an honest record of increasingly disciplined judgment.

\subsection*{2. Recover the Original Reasoning}
Before revising, state what the original work actually did. Identify its claim, evidence or premises, inference, and unstated assumptions. Do not criticize a vague memory of the assignment. Work from the words on the page.

Preserve the original. Revision has meaning only when the reader can see what changed.

\subsection*{3. Diagnose with a Criterion}
``This answer is weak'' is not a diagnosis. A diagnosis names a standard and points to evidence. Examples:

\begin{itemize}[leftmargin=*]
\item the key term shifts between two senses;
\item the premises do not force the conclusion;
\item the source lacks relevant expertise or corroboration;
\item the examples are too narrow for the general claim;
\item the paragraph proves a nearby claim instead of the point at issue;
\item rhetorical force exceeds what the evidence proves.
\end{itemize}

The criterion should fit the artifact. Do not force a syllogism label onto a source-credibility problem merely to sound technical.

\subsection*{4. Make a Substantive Revision}
A substantive revision changes the reasoning. It may clarify a term, narrow a conclusion, expose a hidden assumption, add relevant evidence, remove irrelevant support, repair an inference, qualify a generalization, or distinguish persuasion from proof.

Editing matters, but punctuation alone does not show intellectual growth. Ask: after this change, what claim can the evidence now support that it could not support before?

\subsection*{5. Write the Reflective Caption}
A useful caption answers four questions:

\begin{enumerate}[leftmargin=*,itemsep=0.12em]
\item What did the original artifact try to establish?
\item What precise weakness or limit did I find?
\item Which logical habit guided my revision?
\item What does the revised version now show about my judgment?
\end{enumerate}

Name the change without inventing a heroic transformation. ``I narrowed \textit{all} to \textit{several} because three examples could not support a universal claim'' is stronger than ``I learned to think better.''

\subsection*{6. Build a Revision Trail}
Keep four items together: the original, a short diagnosis, the revised version, and the caption. Label them clearly. A reader should be able to verify the claim about growth without guessing what changed.

For at least one artifact, use \textit{Macbeth} or \textit{Julius Caesar}. Return to the passage rather than relying on plot memory. A later tool---source testing, dispute mapping, rhetoric-vs-proof, or self-deception autopsy---should deepen the original reading rather than merely add vocabulary.

\subsection*{7. Connect Artifacts into a Pattern}
One correction proves that one correction occurred. A portfolio claim becomes stronger when several artifacts show a pattern. Perhaps your conclusions become narrower, your definitions more stable, or your source judgments more careful. Choose a pattern that the artifacts can actually demonstrate.

\subsection*{8. The Six-Step Revision Protocol}
\begin{enumerate}[leftmargin=*,itemsep=0.12em]
\item \textbf{Select} an artifact that can show development.
\item \textbf{Recover} its original claim, support, inference, and assumptions.
\item \textbf{Diagnose} one precise issue using a fitting criterion.
\item \textbf{Revise} the reasoning substantively.
\item \textbf{Caption} the before-and-after change with evidence.
\item \textbf{Connect} the artifact to a larger habit and a new setting.
\end{enumerate}

\section*{Reading Focus}
\begin{enumerate}[leftmargin=*]
\item Why is a portfolio more than a stack of assignments?
\item What makes an artifact useful evidence of growth?
\item Why must the original be preserved?
\item What is the difference between a vague criticism and a criterion-based diagnosis?
\item Name four kinds of substantive revision.
\item What four questions should a reflective caption answer?
\item Copy the six-step revision protocol.
\item Why should at least one revised artifact return to the Shakespeare passage itself?
\end{enumerate}

\section*{Handout Summary}
Write 5--7 sentences explaining how a revision trail can prove intellectual growth. Use the words \textit{artifact}, \textit{criterion}, \textit{substantive revision}, \textit{caption}, and \textit{transfer}.

\section*{Portfolio Preparation}
Before the workshop, bring one early artifact, one later artifact, and one Shakespeare application. For each, mark the sentence or answer area that best records your reasoning. Then prepare these brief notes:
\begin{itemize}[leftmargin=*]
\item a possible pattern of growth that the three artifacts could support;
\item one artifact that needs substantive revision and the criterion that may guide it;
\item one feature of the original that must remain visible in the revision trail;
\item one setting beyond this course where the same logical habit could transfer.
\end{itemize}

\textbf{Working growth claim:} \rule{0.48\columnwidth}{0.4pt}\\[1.1em]
\rule{\columnwidth}{0.4pt}\\[1.1em]
\textbf{First revision move:} \rule{0.48\columnwidth}{0.4pt}
\end{multicols}
\end{document}
